\documentclass[11pt,a4paper]{article}
% Common preamble for the hint sheets (\input from each file)
\usepackage[utf8]{inputenc}
\usepackage[T1]{fontenc}
\usepackage{lmodern}
\usepackage{amsmath,amssymb,amsthm,mathtools}
\usepackage[a4paper,margin=2.3cm]{geometry}
\usepackage{enumitem}
\usepackage{xcolor}
\usepackage[hidelinks]{hyperref}
\renewcommand{\leq}{\leqslant}
\renewcommand{\le}{\leqslant}
\renewcommand{\geq}{\geqslant}
\renewcommand{\ge}{\geqslant}
\setlength{\parindent}{0pt}
\setlength{\parskip}{0.35em}
% \hintstitle{Sheet title}{Class work / Homework}
\newcommand{\hintstitle}[2]{\begin{center}{\LARGE\bfseries Hints: #1\par}\vspace{0.5em}{\large #2\par}\end{center}
\noindent\textit{The hints to each problem are cumulative. Read one, try again, and open the next only if you are still stuck.}\par\medskip\hrule\medskip}
\newcommand{\problem}[1]{\par\bigskip{\large\bfseries Problem #1}\par\nopagebreak\smallskip}
\newcommand{\subpart}[1]{\par\smallskip\textbf{(#1)}\par\nopagebreak}
\newcommand{\hint}[1]{\par\smallskip\noindent\textcolor{black!70}{\textbf{Hint #1.}}\enspace}
\newcommand{\note}[1]{\par\smallskip\noindent\textit{Note.} #1\par}
\providecommand{\lt}{<}
\providecommand{\gt}{>}
\providecommand{\tightlist}{\setlength{\itemsep}{0pt}\setlength{\parskip}{0pt}}
\emergencystretch=3em

\begin{document}
\hintstitle{Abstract algebra -- 1. Finite groups: counting}{Class work}

\problem{1}
\hint{1} For Lagrange: show that two left cosets $xH$, $yH$ are either equal or disjoint, and that $h\mapsto xh$ is a bijection $H\to xH$. The corollaries follow by taking $H=\langle x\rangle$ (whose order is $\operatorname{ord}(x)$), and for Euler by applying $x^{|G|}=e$ in $G=(\mathbb Z/n\mathbb Z)^*$, $|G|=\varphi(n)$.
\hint{2} For the converse, try $G=A_4$ (order $12$). A subgroup of order $6$ would have index $2$, hence be normal, and would contain the square of every element of $A_4$ (look at the image in $G/H\cong\mathbb Z/2$).
\hint{3} Every $3$-cycle is a square: $(abc)=(acb)^2$. So such a subgroup would contain all eight $3$-cycles, which is more than $6$ elements.

\problem{2}
\hint{1} Pair each $x$ with $x^{-1}$. The pairs $\{x,x^{-1}\}$ with $x\ne x^{-1}$ have two elements each; what is left is exactly $\{x: x^2=e\}$. Since $e$ is in that set, $|G|$ even forces at least one more element, an involution.
\hint{2} For McKay: let $X=\{(x_1,\dots,x_p)\in G^p : x_1\cdots x_p=e\}$. The first $p-1$ coordinates are free and the last is determined, so $|X|=|G|^{p-1}\equiv 0\pmod p$.
\hint{3} The cyclic shift maps $X$ to $X$ (conjugate $x_1\cdots x_p=e$ by $x_1$), and its $p$-th power is the identity, so orbits have size $1$ or $p$ ($p$ is prime). Hence the number of fixed points, the constant tuples $(x,\dots,x)$ with $x^p=e$, is $\equiv |X|\equiv 0\pmod p$. As $(e,\dots,e)$ is one of them, there are at least $p$, so some $x\ne e$ has $x^p=e$.

\problem{3}
\hint{1} The $n$ elements of order $p$ together with $e$ are exactly the solutions of $x^p=e$. So you need: the number of solutions of $x^p=e$ is divisible by $p$ whenever $n>0$.
\hint{2} If $n>0$, then $p$ divides $|G|$ by Lagrange. Now rerun McKay's argument from Problem 2: the fixed points of the cyclic shift on $X=\{(x_1,\dots,x_p):x_1\cdots x_p=e\}$ are the tuples $(x,\dots,x)$ with $x^p=e$, and their number is $\equiv|G|^{p-1}\equiv0\pmod p$. Hence $p\mid n+1$.

\problem{4}
\hint{1} A unique subgroup of order $p$ is normal, because every conjugate $gHg^{-1}$ is again a subgroup of order $p$. So conjugation gives a homomorphism $G\to\operatorname{Aut}(H)$.
\hint{2} $H\cong\mathbb Z/p$, so $|\operatorname{Aut}(H)|=p-1$. The image of $G$ has order dividing both $n$ and $p-1$.
\hint{3} Every prime divisor of $n$ is $\ge p$, so $\gcd(n,p-1)=1$. The action is therefore trivial, i.e.\ every $g$ commutes with every element of $H$.

\problem{5}
\hint{1} If all orders multiply to a power of $2$, every element has $2$-power order, so $G$ is a $2$-group: $G\cong\bigoplus_i\mathbb Z/2^{a_i}$ with $|G|=2^n$ and $r$ cyclic summands. The product of orders is $2^E$. Let $N(j)=\#\{x:x^{2^j}=e\}$, which is a power of $2$. Write $E$ in terms of the $N(j)$.
\hint{2} With $M$ the largest $a_i$, $E=\sum_{j=1}^{M}j\,(N(j)-N(j-1))=M\cdot 2^n-\sum_{j=0}^{M-1}N(j)$. Here $N(0)=1$, and $N(1)=2^r$ divides every $N(j)$ with $j\ge1$ as well as $2^n$. Hence $E\equiv-1\pmod{2^r}$.
\hint{3} $E+1=2010$ is not divisible by $4$, so $r\le1$ and $G$ is cyclic, $G=\mathbb Z/2^n$. Then $E=\sum_{j=1}^n j\,2^{j-1}=(n-1)2^n+1$, and $(n-1)2^n=2008$ has no solution. The answer is \emph{no}.

\end{document}
